% !TEX root = /home/eb/.vscode/extensions/qftrules.programme-colle/tmp/Exercice.tex

%%%%%%%%%%%%%%%%%%%%%%%%%%%%%%%%%%%
\begin{exo}[PC][2][TD]{Expérience de Davisson-Germer}
  \Notions{\item Relation de de Broglie \task Réseau de diffraction}
  En 1927, C. \textsc{Davisson} et L. \textsc{Germer} bombardent une cible de nickel monocristallin par des électrons dont ils maîtrisent la vitesse grâce à un canon à électrons utilisant une tension électrique ajustable. Un détecteur mobile permet ensuite d'étudier la dépendance angulaire du faisceau diffracté. On observe que les électrons sont diffractés par la cible principalement dans une direction privilégiée faisant un angle $\theta \simeq \SI{50}{^\circ}$ dans la direction opposée au faisceau incident. Les électrons du faisceau incident , de charge $-e = \SI{-1.60e-19}{C}$ et de masse $m=\SI{9.11e-31}{kg}$, sont accélérés entre les armatures d'un condensateur sous une différence de potentiel $U = \SI{54}{V}$. On donne la constante de Planck $h=\SI{1.05e-34}{J.s}$.

  \begin{quest}
    \item	Montrer qu'une particule de charge $q$ soumise à un potentiel électrostatique $V$ possède une énergie potentielle d'interaction $\Epot = qV$.
    \sol{
      L'énergie potentielle d'une particule de charge $q$ dans un potentiel électrostatique $V$ est donnée par la relation
      % $$
      \eqbox{
        \Epot = qV.
        }
      % $$
      En effet, la force électrostatique est donnée par la relation $\F = q \E = -q \grad V = - \grad \Epot$.
    }
    \item En déduire la vitesse $v$ d'un électron en sortie du condensateur, sachant que les électrons ont une vitesse initiale nulle. Déterminer la longueur d'onde de de Broglie $\lambda$ de l'électron. \AN. À quel domaine du spectre électromagnétique correspondrait cette longueur d'onde ?
    \sol{
    La conservation de l'énergie mécanique donne 
    $$
    0 -e V_1 = \frac{1}{2}mv^2 -eV_2 \Rar \frac{1}{2}mv^2 = eU,
    $$
    soit 
    \eqbox{
      v = \sqrt{\frac{2eU}{m}} = \SI{1.4e7}{m.s^{-1}}.
    }
    La relation de de Broglie s'écrit
    $$
    \lambda = \frac{h}{p} = \frac{h}{mv},
    $$
    soit
    \eqbox{
      \lambda  = \frac{h}{\sqrt{2meU}} = \SI{0.24}{nm}.
    }
    }
  \end{quest}
    
  On admet que la cible de nickel est un réseau cristallin de type cubique à faces centrées, dont les plans agissent comme un réseau de diffraction de pas $a$.
  \begin{quest}
    \item	Rappeler la formule des réseaux, appelé relation de Bragg dans le contexte de la diffraction sur les solides cristallins. Déterminer $a$, assimilée à la distance interatomique dans le cristal de nickel.
    \sol{
    La formule des réseaux donne 
    $$
    \sin i' - \sin i = p\frac{\lambda}{a}.
    $$
    Ici l'angle d'incidence vaut $i' = 0$ et l'angle de diffraction vaut $i = \theta$ (on suppose  qu'il s'agit ici d'un réseau en réflexion et non en transmission, mais cela conduit à la même formule avec les mêmes angles). En supposant que l'ordre d'interférence observé est $p=1$, on a
    \eqbox{
    a = \frac{\lambda}{\sin \theta} = \SI{0.31}{nm}.
    }
    }
  \end{quest}
\end{exo}


%%%%%%%%%%%%%%%%%%%%%%%%%%%%%%%%%%%
\begin{exo}[PC][1][TD]{Modèle planétaire de l'atome d'hydrogène}
  \Notions{\item Inégalite de Heisenberg \task Modèle atomique de Bohr}
Le modèle le plus simple de l'atome d'hydrogène consiste en un modèle planétaire, où le noyau est supposé fixe et autour duquel \guill{gravite} un électron sur une orbite circulaire stable de rayon $r$ et de quantité de mouvement $p$. Sur la première orbite de Bohr, le rayon vaut $ a_0=\SI{0.5e-10}{m}$ et la quantité de mouvement $p_0=\hbar/a_0$.
% (sur laquelle il ne rayonne aucune énergie).
% Sur la première orbite de, N. Bohr postula que la quantité de mouvement associée à l'électron vaut $p_0=\hbar/a_0$. 
L'objectif de cet exercice est de discuter de la pertinence de ce modèle du point de vue de la mécanique quantique.
\begin{quest}
\item Évaluer $p_0$. Pour pouvoir définir une trajectoire classique, quelle limitation doit-on imposer aux indéterminations $\Delta r$ et $\Delta p$ sur le rayon et la quantité de mouvement pour la première orbite de Bohr ?
 
\sol{
Pour que la trajectoire classique soit bien définie, on doit avoir $\Delta r \ll a_0$ et $\Delta p \ll p_0$. On devrait donc avoir
% $$
\eqbox{
  \Delta r \Delta p \ll a_0 p_0 = \hbar.
}
% $$
}
\item Montrer que ces limitations sont incompatibles avec l'inégalité de Heisenberg, et justifier la nécessité d'un traitement quantique de l'atome.
\sol{
L'inégalité de Heisenberg impose
% $$
\eqbox{
  \Delta r \Delta p \geq \frac{\hbar}{2},
}
ce qui est incompatible avec la condition précédente. Un traitement quantique est donc nécessaire.
% $$
}
% Commenter.
\end{quest}
\end{exo}
%%%%%%%%%%%%%%%%%%%%%%%%%%%%%%%%%%%%%%%%%%%%%%%%%%%%%%%%%%%%%%%%%%

%%%%%%%%%%%%%%%%%%%%%%%%%%%%%%%%%%%
\begin{exo}[PC][1][appli]{Calcul de grandeurs ondulatoires}
  \Notions{\item Relation de de Broglie \task Relation de Planck-Einstein}
  \hauteurLargeurCadreReponse		{8mm}{7cm}
Calculer numériquement la longueur de de Broglie $\lambda$ et la fréquence $\nu$ associée à l’onde quantique dans les situations suivantes et commenter.

  \begin{enonce}
  	Atome de néon de masse $m=\SI{3.35e-26}{kg}$ de vitesse $v=\SI{3}{m.s^{-1}}$ dans une expérience de diffraction de fentes d'Young.\\

  \end{enonce}

  % \reponse{$r \left(\cos(\theta) \ex+\sin(\theta) \ey\right)$}

  \begin{corrige}

  \end{corrige}
  
  \begin{enonce}
    Atome de fer de masse $m=\SI{9.27e-26}{kg}$ dans un réseau cristallin de température $T=\SI{300}{K}$.\\ 

  \end{enonce}

\end{exo}
%%%%%%%%%%%%%%%%%%%%%%%%%%%%%%%%%%%

%%%%%%%%%%%%%%%%%%%%%%%%%%%%%%%%%%%
\begin{exo}[PC][3][colle]{Supraconducteur et jonction Josephson}
  \Notions{\item Équation de Schrödinger \task État non stationnaire}
  Les supraconducteurs sont des matériaux qui présentent une absence totale de résistivité électrique en dessous d'une certaine température critique $T_c$. Le premier supraconducteur découvert est le mercure par K. Onnes en 1904, et présente une température critique de $\SI{4.2}{K}$.
  La supraconductivité apparaît lorsque les électrons de conduction du matériau forment des états liés appelés paires de Cooper.
  % \begin{fig}[label]
  %   \centering
  %   \begin{subfig}[label][0.43]
  %     \includegraphics[height=0.64\columnwidth]{supra.jpg}
  %   \end{subfig}\hfill
  %   \begin{subfig}[label][0.43]
  %     \includegraphics[height=0.64\columnwidth]{josephson.jpg}
  %   \end{subfig}
  %   \capt{
  %     (a) Photographie d’un matériau supraconducteur. (b) Image d’une jonction Josephson par microscopie électronique à balayage [Institute for Quantum Computing, University of Waterloo].
  %   }
  % \end{fig}
  \begin{quest}
    \item	On décrit l’assemblée de paires de Cooper par un fluide quantique de fonction d’onde $\psi(x,t) = \sqrt{n}\,e^{i\alpha(t)}$ où $n$ représente la densité volumique de paires de Cooper, supposée constante, et $\alpha(t)$ désigne la phase de la fonction d’onde, supposée uniforme. Justifier l’interprétation physique de $n$.
    %  proposée à partir d’une analyse dimensionnelle. 
    % Que dire de la densité de probabilité $\rho(x,t)$ ?
    \item On rappelle que l’énergie potentielle d’un électron dans un potentiel électrostatique $V$ s’écrit $E=-eV$, où $e$ est la charge élémentaire. Écrire l’équation de Schrödinger pour la fonction d’onde $\psi(x,t)$, et en déduire l’équation différentielle sur $\alpha(t)$. Que vaut $\alpha(t)$ pour un potentiel électrostatique $V$ constant ? 
    \sol{
      L’équation de Schrödinger s’écrit
      $$
      i\hbar\dd{\psi}{t} = -\frac{\hbar^2}{2m}\dd{\psi}{x} - eV\psi \Rar -\hbar \sqrt{n} \dd{\alpha}{t}e^{i\alpha(t)} = -e V \sqrt{n} e^{i\alpha(t)},
      $$
      et donc 
      \eqbox{
        \dd{\alpha}{t} = \frac{eV}{\hbar}.
      }
      Si le potentiel $V$ est constant, on a 
      $$
      \alpha(t) = \frac{eV}{\hbar}t.
      $$
    }
  \end{quest}
  On considère un composant formé de deux échantillons supraconducteurs séparés par une fine couche isolante ; ce composant est appelé jonction Josephson. Les paires de Cooper dans les supraconducteurs $1$ et $2$ sont décrits par les fonctions d’onde 
  $$\psi_1(t) = \textstyle\sqrt{n_1(t)}\,e^{i\alpha_1(t)}, \qquad\text{et}\qquad \psi_2(t) = \textstyle\sqrt{n_2(t)}\,e^{i\alpha_2(t)},$$ respectivement, où les densités volumiques $n_1(t)$ et $n_2(t)$ dépendent cette fois-ci du temps. On admet que l’effet tunnel au niveau de la jonction introduit un couplage entre les fonctions d’onde, qui vérifient alors les équations  de Schrödinger
    % $$

\begin{equation}
  i\hbar \frac{\mathrm{d} \psi_1}{\mathrm{~d} t}=K \psi_2, \quad \text { et } \quad i\hbar \frac{\mathrm{d} \psi_2}{\mathrm{~d} t}=K \psi_1.
\end{equation}  
    % $$
    
    \begin{quest}
    \item 
    Montrer que la variation de population des paires de Cooper dans le supraconducteur $1$ vérifie
    % $$
    \begin{equation}
      \dd{n_1}{t} = \frac{2K}{\hbar}\sqrt{n_1n_2} \,\sin (\alpha_2-\alpha_1).
    \end{equation}
    % $$

    \sol{
    On commence par calculer 
    $$
    \dd{n_1}{t} = \dd{}{t}(\psi_1^* \psi_1) = \dd{\psi_1^*}{t} \psi_1 + \psi_1^* \dd{\psi_1}{t} = \frac{iK}{\hbar}\psi_2^* \psi_1 - \frac{iK}{\hbar}\psi_2 \psi_1^* = \frac{iK}{\hbar} \times (-2i) \im(\psi_2 \psi_1^*),
    $$
    d'où le résultat,
    \eqbox{
      \dd{n_1}{t} = \frac{2K}{\hbar}\sqrt{n_1n_2} \,\sin (\alpha_2-\alpha_1).
    }
    }
    \item Sachant que les paires de Cooper possède une charge électrique $-2e$, déterminer l’expression de l’intensité $I$ du courant qui traverse la jonction Josephson en fonction de $e$, $K$, $n_1$, $n_2$, $\alpha_2$, $\alpha_1$, et du volume $v$ du supraconducteur $1$. Exprimer $\alpha_2 - \alpha_1$ lorsque les supraconducteurs sont portés à une tension $U=V_1-V_2$ constante. En déduire la caractéristique d’une jonction Josephson,
    % $$
    \begin{equation}
      I(U,t) = I_0 \sin\left(\pi\frac{Ut}{\Phi_0}\right),
    \end{equation}
    % $$
    où l’on donnera les expressions de $I_0$ et $\Phi_0$. Quelle est la dimension et la valeur numérique de $\Phi_0$ ?
    \sol{
    L’intensité du courant est relié au nombre de particule $q_1$ qui traverse la jonction de $2$ vers $1$ par 
    $$
    I = \dd{q_1}{t} = \dd{}{t}(-2e n_1 v) = -2e v \dd{n_1}{t}  = -\frac{2Kev}{\hbar}\sqrt{n_1n_2} \,\sin (\alpha_2-\alpha_1).
    $$
    Lorsque la jonction est soumise à une tension $U=V_1-V_2$ constante, on a 
    $$
    \alpha_2 - \alpha_1 = -\frac{eU}{\hbar}t,
    $$
    d'où 
    \eqbox{
    I = \frac{2Ke}{\hbar}v\sqrt{n_1n_2} \,\sin \left(\frac{eUt}{\hbar}\right).
    }
    En supposant que les densités particulaires $n_1$ et $n_2$ varient peu, la grandeur 
    $$
    I_0 = \frac{2Ke}{\hbar}v\sqrt{n_1n_2}
    $$
    est une constante, homogène à une intensité de courant. La constante  $\Phi_0$ est définie par
    $$
    \Phi_0 = \frac{h}{2e} = \SI{3.2e-15}{Wb}.
    $$
    Il s'agit du quantum de flux magnétique.
    }
    
  \end{quest}
\end{exo}
%%%%%%%%%%%%%%%%%%%%%%%%%%%%%%%%%%%


%%%%%%%%%%%%%%%%%%%%%%%%%%%%%%%%%%%%%%%%%%%%%%%%%%%%%%%%%%%%%%%%%%
\begin{exo}[PC][1][TD]{Analogie entre optique et mécanique quantique}
  \Notions{\item Interférences \task État stationnaire}
  On considère une onde lumineuse harmonique de vibration scalaire $s(x,t)$ issue de la superposition de deux ondes cohérentes harmoniques $s_1(x,t)$ et $s_2(x,t)$ de même pulsation $\omega$.
\begin{quest}
\item  Définir l'intensité lumineuse $I(x)$ et démontrer la formule des interférences.
\sol{
Superposition : $s(x,t) = s_1(x,t) +  s_2(x,t) = A \cos(\omega t - k x) + A \cos(\omega t - k x + \vphi)  = 2A \cos(\vphi/2)\cos(\omega t - k x)$. Définition de l'intensité lumineuse :
$$
I(x) = \moy{s(x,t)^2 } = 4A^2\cos(\vphi/2)^2 \moy{\cos(\omega t - k x)^2} = 2A^2\cos(\vphi/2)^2,
$$
soit 
\eqbox{
  A^2\left[ 1 + \cos(\vphi)\right].
}
}
\end{quest}
On considère maintenant un électron décrit par une fonction d'onde stationnaire $\psi(x,t)$ issue de la superposition des deux fonctions d'ondes $\psi_1(x,t) = A e^{-iEt/\hbar}$ et $\psi_2(x,t) = A e^{i\vphi(x)} e^{-iEt/\hbar}$.
\begin{quest}
\item  Rappeler la définition de la densité de probabilité $\rho(x,t)$ et la calculer. Comparer à la formule des interférences en optique. Quelles sont les similarités et les différences ?
\sol{
Superposition : $\psi(x,t) = \psi_1(x,t) +  \psi_2(x,t) = A e^{-iEt/\hbar} (1+e^{i\vphi})$. Définition de la densité de probabilité :
$$
\rho(x,t) = |s(x,t)|^2 = A^2 (1+e^{i\vphi})(1+e^{-i\vphi}),
$$
soit 
\eqbox{
 \rho(x,t) =  2A^2\left[1+ \cos(\vphi) \right].
}
\udl{similarités} : le principe le superposition et la forme des interférences s'appliquent de la même façon.

\udl{différences} : $\psi$ est intrinsèquement complexe alors que $s$ est réel. $I$ s'obtient par une moyenne temporelle en raison du temps de réponse du photorécepteur, alors que $\rho$ peut dépendre du temps.
}
\end{quest}
\end{exo}
%%%%%%%%%%%%%%%%%%%%%%%%%%%%%%%%%%%%%%%%%%%%%%%%%%%%%%%%%%%%%%%%%%


%%%%%%%%%%%%%%%%%%%%%%%%%%%%%%%%%%%%%%%%%%%%%%%%%%%%%%%%%%%%%%%%%%
\begin{exo}[PC][1][TD]{Interférences quantiques}
  \Notions{\item Interférences \task Fentes de Young}
On considère une expérience d'interférences de types fentes de Young avec un faisceau d'électrons. Le système contient 3 fentes qui peuvent être fermées ou ouvertes, et on détecte les particules en un point $M$ d'un écran situé à grande distance des fentes. On admet que:
\begin{itemize}
\item lorsque seule la fente 1 est ouverte, l'amplitude de probabilité en $M$ vaut $\psi_1(M)=\frac{1}{\sqrt{3}}$;
\item lorsque seule la fente 2 est ouverte, l'amplitude de probabilité en $M$ vaut $\psi_2(M)=\frac{i}{\sqrt{2}}$;
\item lorsque seule la fente 3 est ouverte, l'amplitude de probabilité en $M$ vaut $\psi_3(M)=\frac{-1}{\sqrt{6}}$.
\end{itemize}
\begin{quest}
\item Déterminer la densité de probabilité de détection d'une particule en $M$ lorsque: a) seule la fente 2 est ouverte; b) les fentes 1 et 2 sont ouvertes; c) toutes les fentes sont ouvertes.
\sol{
La densité de probabilité vaut
$\rho(M) = |\psi(M)|^2$.
\begin{itemize}
\item lorsque seule la fente 2 est ouverte, l'amplitude de probabilité en $M$ vaut $\psi_2(M)=\frac{i}{\sqrt{2}}$, et donc $\rho(M) = |\psi_2(M)|^2 = \frac{1}{2}$;
\item lorsque seule les fentes 1 et  2 sont ouvertes, l'amplitude de probabilité en $M$ vaut par superposition, $\psi(M)=\psi_1(M)+\psi_2(M)=\frac{1}{\sqrt{3}}+\frac{i}{\sqrt{2}}$, et donc
$\rho(M) = |\psi(M)|^2 = \frac{1}{3}+\frac{1}{2} = \frac{5}{6}$ ;
\item lorsque toutes les fentes sont ouvertes, l'amplitude de probabilité en $M$ vaut $\psi(M)=\psi_1(M)+\psi_2(M)+\psi_3(M)=\frac{1}{\sqrt{3}}+\frac{i}{\sqrt{2}}+\frac{-1}{\sqrt{6}} $, et donc $\rho(M) = |\psi(M)|^2 = \frac{1}{6}+\frac{1}{2} = (\frac{1}{\sqrt{3}}-\frac{1}{\sqrt{6}})^2 +\frac{1}{2} \simeq \SI{0.53}{}.$
\end{itemize}
}
\item Commenter sur le fait qu'en ouvrant la fente 3, on diminue la probabilité de détecter une particule en $M$ par rapport au cas où les fentes 1 et 2 sont ouvertes.  
% Existe-t-il un analogue classique de ce résultat ?
\sol{
On peut comparer cela aux interférences en optique. En représentation de Fresnel, on voit qu'on somme des nombres complexes, i.e. des vecteurs du plan, qui sont en opposition de phase. On peut aussi comparer à l'effet inverse observé lorsqu'on augmente la luminosité en sortie d'un système polariseur-analyseur croisé en ajoutant un troisième polariseur à $\SI{45}{\degres}$ intercalés entre les deux, même ce ne sont pas les mêmes relations mathématiques qui interviennent.
}
\end{quest}
\end{exo}
%%%%%%%%%%%%%%%%%%%%%%%%%%%%%%%%%%%%%%%%%%%%%%%%%%%%%%%%%%%%%%%%%%

%%%%%%%%%%%%%%%%%%%%%%%%%%%%%%%%%%%%%%%%%%%%%%%%%%%%%%%%%%%%%%%%%%
\begin{exo}[PC][2][TD]{Oscillateur harmonique quantique}
  \Notions{\item État stationnaire \task Normalisation \task Inégalité de Heisenberg}
Soit un atome de carbone de masse $m=\SI{2.0e-26}{kg}$ piégé dans le réseau cristallin d’un échantillon de diamant. En raison de l’interaction avec les atomes voisins, l'atome de carbone évolue dans un potentiel harmonique, de la forme
\begin{equation}
  V(x) = \textstyle\frac{1}{2}m\omega^2 x^2,
\end{equation}
de pulsation $\omega = \SI{8.7e11}{rad.s^{-1}}$. On cherche des états stationnaires de l'atome de carbone.
\begin{quest}
\item Rappeler la forme d'une fonction d'onde stationnaire  $\psi(x,t)$ d'une particule d'énergie $E$. Écrire l'équation de Schrödinger indépendante du temps qui porte sur la partie spatiale $\vphi(x)$ de la fonction d'onde.
\sol{
On a $\psi(x,t)=\varphi(x)\,\e^{-iEt/\hbar}$. L'équation de Schrödinger indépendante du temps s'écrit
% $$
\eqbox{
  -\frac{\hbar^2}{2m}\frac{\d^2\varphi}{\d x^2}+ \frac{1}{2}m\omega^2x^2\varphi =E\varphi.
  }
% $$
}
\item Vérifier qu'une fonction de la forme $\varphi(x)=A \,\e^{-x^2/a^2}$, avec $A$ et $a$ des constantes, satisfait à l'équation précédente. En déduire l'expression de la distance caractéristique $a$ et de l'énergie $E$ en fonction de $\hbar$, $m$ et~$\omega$. \emph{Application numérique}.

\sol{
  On a
  \begin{align*}
    -\frac{\hbar^2}{2m}\frac{\d^2\varphi}{\d x^2} = -\frac{\hbar^2}{2m}\dd{}{x}\left( \dd{}{x} A \,\e^{-\frac{x^2}{a^2}}\right) & = -A\frac{\hbar^2}{2m}\dd{}{x}\left( -\frac{2x}{a^2} \,\e^{-\frac{x^2}{a^2}}\right) \\
    & = A\frac{\hbar^2}{ma^2}\left(1 - 2\frac{x^2}{a^2}\right) \,\e^{-\frac{x^2}{a^2}} \\
    & = \frac{\hbar^2}{ma^2} \varphi - \frac{2\hbar^2}{ma^4} x^2 \vphi.
  \end{align*}
  L'équation est vérifiée si  $\frac{2\hbar^2}{ma^4} = \frac{1}{2}m\omega^2$, soit
  \eqbox{a = \sqrt{\frac{2\hbar}{m\omega}},}
  et
  \eqbox{E = \frac{\hbar^2}{ma^2} = \frac{\hbar \omega}{2}.}
  }
\end{quest}
  \emph{Il s’agit du mode fondamental de l’oscillateur harmonique quantique. Il est possible de calculer l'énergie du mode d'ordre $n\in \N$, dont l'expression est $E_n = \hbar \omega \textstyle(n+\frac{1}{2})$.}
\begin{quest}
\item
Déterminer la constante de normalisation $A$. \emph{Application numérique}.
\sol{
  Supposons $A \in \R^+$ sans perte de généralité. On a
  $$
\int_{-\infty}^{+\infty} |\vphi(x)|^2 \d x= A^2 \int_{-\infty}^{+\infty} \e^{-2\frac{x^2}{a^2}} \d x = A^2 \sqrt{\frac{\pi a^2}{2}} = 1 ,
$$
donc \eqbox{A = \sqrt{\frac{1}{a}\sqrt{\frac{2}{\pi}}}.}
}
\item Représenter le graphe de la densité de probabilité de présence  $\rho(x)$. En déduire sans calcul explicite la valeur moyenne de la position $\avg{x}$. Commentaire ? Calculer l'indétermination $\Delta x$ sur la position, et en déduire une inégalité sur l'indétermination de la quantité de mouvement $\Delta p_x$. \emph{Application numérique}.
\sol{
La densité de probabilité de présence se calcule comme
$$
\rho(x) = |\vphi(x)|^2 = A^2  \e^{-2\frac{x^2}{a^2}} = \sqrt{\frac{2}{\pi a^2}} \e^{-2\frac{x^2}{a^2}}.
$$
Il s'agit d'une fonction paire, donc centrée autour de $\avg{x} = 0$. L'écart-type $\Delta x$ sur la position vérifie
$$
\Delta x^2 = \avg{x^2} - \avg{x}^2 = \avg{x^2} =  A^2 \int_{-\infty}^{+\infty} x^2 \e^{-2\frac{x^2}{a^2}} =  \sqrt{\frac{2}{\pi a^2}} \times \frac{1}{2} \sqrt{\frac{\pi a^6}{8}} = \sqrt{\frac{a^4}{16}} = \frac{a^2}{4},
$$
soit
\eqbox{\Delta x = \frac{a}{2}.}
L'indétermination $\Delta p_x$ s'obtient par l'inégalité de Heisenberg, soit
\eqbox{\Delta p_x \geq \frac{\hbar}{2\Delta x} = \frac{\hbar}{a}.}
}
\sol{
À titre informatif, on donne ci-dessous les premiers états stationnaires (la fonction d'onde est réelle) de l'oscillateur harmonique quantique, calculés sur python par résolution numérique de l'équation de Schrödinger.
\begin{fig}[label]
  \centering
  \includegraphics[width=0.7\columnwidth]{schrodinger_Quadraticpotential}
  %\capt{}
\end{fig}
}
% \item Déterminer l'expression de l'énergie $E$ et de la distance $a$ en fonction de $\hbar,\omega$ et $m$.
\end{quest}
\textsc{Données} : $\int_{-\infty}^{\infty}\e^{-\alpha u^2}\d u  =\sqrt{\frac{\pi}{\alpha}}$, $\int_{-\infty}^{\infty}u^2\e^{-\alpha u^2}\d u  =\frac{1}{2} \sqrt{\frac{\pi}{\alpha^3}}$.
\end{exo}
%%%%%%%%%%%%%%%%%%%%%%%%%%%%%%%%%%%%%%%%%%%%%%%%%%%%%%%%%%%%%%%%%%


%%%%%%%%%%%%%%%%%%%%%%%%%%%%%%%%%%%%%%%%%%%%%%%%%%%%%%%%%%%%%%%%%%
\begin{exo}[PC][1][colle]{Fluctuations d’un oscillateur harmonique quantique}
  \Notions{\item Équation de Schrödinger indépendante du temps \task Relation d'indétermination de Heisenberg}
Une particule quantique de masse $m$ évolue dans un potentiel harmonique $V(x)=\frac{1}{2}m\omega^2 x^2$, où $\omega$ est homogène à une pulsation. On admet que la position moyenne et la quantité de mouvement moyenne sont nulles en raison de la parité du potentiel. On note $\Delta p$ et $\Delta x$ les indéterminations sur la quantité de mouvement et la position de la particule.
\begin{quest}
\item À partir de l'inégalité de Heisenberg, montrer que l'énergie de la particule $E=E_c + V$  est bornée inférieurement par une fonction $E_{\rm inf}(\Delta x)$ à déterminer.
\sol{
  Nous utilisons ici un argumert semi-classique.
  L'énergie totale classique de la particule vaut 
  $$
  E = E_c + V = \frac{p^2}{2m} + \frac{1}{2}m\omega^2 x^2.
  $$
  En mécanique quantique, les indéterminations sur $x$ et $p$ sont liées par l'inégalité d'Heisenberg, 
  $$
  \Delta x \Delta p \geq \frac{\hbar}{2}.
  $$
  Par définition, $\Delta x^2 = \moy{x^2} - \moy{x}^2$ et $\Delta p = \moy{p^2} - \moy{p}^2$. Comme $\moy{x} = \moy{p} = 0$, on a $\Delta x = \sqrt{\moy{x^2}}$ et $\Delta p = \sqrt{\moy{p^2}}$. L'ordre de grandeur typique de la position et de la quantité de mouvement valent donc $x \sim \sqrt{\moy{x^2}} = \Delta x$ et $p \sim \sqrt{\moy{p^2}} = \Delta p \geq \frac{\hbar}{2\Delta x}$. On en déduit que $E \geq E_{\rm inf}(\Delta x)$ avec 
  \eqbox{
    E_{\rm inf}(\Delta x) = \frac{\hbar^2}{8m\Delta x^2} + \frac{1}{2}m\omega^2 \Delta x^2.
  }
}
\item En déduire l'indétermination sur la position $\Delta x_0$ en fonction de $m$, $\hbar$ et $\omega$ lorsque la particule est dans son état fondamental, d'énergie minimale $E_{0}$ à exprimer en fonction de $\hbar$ et $\omega$.
\sol{
La fonction $E_{\rm inf}(\Delta x)$ admet un minimum lorsque 
$$
\dd{E_{\rm min}}{\Delta x}(\Delta x_0) = 0 \Rar -\frac{\hbar^2}{4m} \frac{1}{\Delta x_0^3} + m\omega^2 \Delta x_0 =0,
$$
\eqbox{
  \Delta x_0 = \sqrt{\frac{\hbar}{2m\omega}}.
}
On trouve ensuite $E_0 \geq E_{\rm inf}(\Delta x_0)$. En supposant que la borne inférieure est atteinte, autrement dit que l'inégalité de Heisenberg est saturée pour l'état fondamental,
 on trouve
\eqbox{
E_0 = \frac{1}{2} \hbar \omega.
}
Cette valeur correspond à la valeur exacte de l'énergie de l'état fondamental de l'oscillateur harmonique quantique. En réalité, il s'agit d'un coup de chance, l'argument précédent ne donne qu'une borne inférieure sur l'énergie en ordre de grandeur.
}
\item L’agitation thermique engendre des fluctuations d'énergie $k_{\rm B} T$. Déterminer la température critique~$T_{\rm{c}}$ en dessous de laquelle les fluctuations quantiques dominent les fluctuations thermiques. 
\sol{
  Les fluctuations quantique dominent si $E > k_B T$ pour tout état d'énergie $E$. L'état le plus contraignant est l'état fondamental, d'énergie minimale $E_0$. La température critique $T_c$ qui sépare les deux régimes vérifie donc $E_0 = k_B T_c$, soit
  \eqbox{
  T_c = \frac{\hbar \omega}{2k_B}.
  }
% Les fluctuations quantiques dominent si $\Delta x > \Delta x_{\rm th}$, soit
% $$
% \sqrt{\frac{\hbar}{m\omega}} > \sqrt{\frac{k_B T}{m\omega^2}} \Rar T < T_{\rm c} = \frac{\hbar\omega}{k_B}.
% $$
}
\item Estimer $T_{\rm{c}}$ pour un oscillateur mécanique masse-ressort classique de laboratoire et commenter. En 2010, une équipe de l'université de Santa Barbara a atteint le régime quantique en amenant un micro-résonateur piézo-électrique de fréquence $f=\SI{6}{GHz}$ à une température de $\SI{25}{mK}$. Commenter ce choix de température. Estimer les fluctuations quantiques de position.
\sol{
On trouve ici \fbox{$T_{\rm c} = \SI{0.32}{K}$}. Comme $T<T_c$, les phénomènes quantiques sont bien visibles. L'indétermination sur la position dans le mode fondamental vaut ici $\Delta x_0 = \SI{0.1}{pm}$.
}
\end{quest}
\end{exo}
%%%%%%%%%%%%%%%%%%%%%%%%%%%%%%%%%%%%%%%%%%%%%%%%%%%%%%%%%%%%%%%%%%%%%%%%%%%%%%%%%%%

%%%%%%%%%%%%%%%%%%%%%%%%%%%%%%%%%%%%%%%%%%%%%%%%%%%%%%%%%%%%%%%%%%%%%%%%%%%%%%%%%%%
\begin{exo}[PC][2]{Propagation d’un paquet d'onde gaussien}
  \Notions{\item Particule libre \task Paquet d'onde \task Inégalité de Heisenberg \task Normalisation}
  On souhaite décrire de façon réaliste la propagation d’une particule quantique libre.
\begin{quest}
\item Rappeler l'expression des fonctions d'onde stationnaires solutions de l'équation de Schrödinger pour une particule libre. Pourquoi une OPPH $\psi(x,t) = Ae^{i(kx-\omega t)}$ ne peut pas décrire une particule réelle ?
\sol{
  La fonction d'onde stationnaire s'écrit 
  \eqbox{
    \psi(x,t) = A e^{i(kx-\omega t)},
  }
  où la pulsation $\omega$ et le vecteur d'onde $k$ sont reliés par la relation de dispersion 
  $$
  \omega(k) = \frac{\hbar k^2}{2m}.
  $$
  Une OPPH ne peut pas décrire une particule réelle car elle n'est pas normalisable.
}
\end{quest}
On construit un paquet d’onde par superposition d’un continuum d’OPPH, de la forme 
\begin{equation*}
\psi(x,t)=\int_{-\infty}^{+\infty}A(k)\, \e^{i(k x - \omega t)}\d k,
\end{equation*}
où l’on prend une fonction réelle gaussienne
$A(k) = A_0 \, e^{-k^2/4\Delta k^2}$ pour l’amplitude spectrale.
% On considère un paquet d'onde gaussien, qui dépend de deux constantes réelles  $A_0$ et $\Delta k$ telles que
% \begin{equation*}
% A(k) = A_0 \exp\left[\, -\frac{k^2}{4\Delta k^2}\right].
% \end{equation*}

\begin{quest}
\item Que représente la grandeur $\Delta k$ ?  Exprimer la fonction d'onde initiale $\psi(x,t=0)$ puis la densité de probabilité correspondante $\rho(x,t=0)$.
% en se servant de la formule de l’intégrale gaussienne 
% $$
% \int_{-\infty}^{+\infty} \e^{-ak^2+bk}\d k = \sqrt{\frac{\pi}{a}}\,\e^{\frac{b^2}{4a}}.
% $$
En déduire la position moyenne $\moy{x}$ et la largeur spatiale $\Delta x$ du paquet d’onde à l'instant initial $t=0$. Vérifier que l’inégalité de Heisenberg $\Delta x \Delta p \geq \frac{\hbar}{2}$ est satisfaite.
\sol{
  
  On a
  % $$
  \begin{align*}
    \psi(x,t = 0) & = \int_{-\infty}^{+\infty} A_0 \exp\left[\, -\frac{(k-k_0)^2}{4\Delta k^2}\right] \e^{ikx}\d k \\ 
    & = A_0 e^{ i\frac{k_0 x}{4 \Delta k^2}} \int_{-\infty}^{+\infty} \exp\left[\, -\frac{k^2}{4\Delta k^2} + ikx\right]\d k,
  \end{align*}
  soit 
  \eqbox{
    \psi(x,t = 0) = A_0\, e^{ i \frac{k_0 x}{4 \Delta k^2}} \sqrt{4\pi\Delta k^2} \exp\left[\, - \Delta k^2 x^2 \right].
  }
  On en déduit la densité de probabilité
  \eqbox{
    \rho(x,t=0) = |\psi(x,t=0)|^2 = |A_0|^2 4\pi \Delta k^2 \exp\left[\, - 2 \Delta k^2 x^2 \right].
  }
  La densité de probabilité est une gaussienne centrée autour de $\moy{x}=0$ et d'écart-type $\Delta x = \frac{1}{2\Delta k}$. On remarque que $\Delta x \Delta p = \Delta x \hbar \Delta k = \frac{\hbar}{2}$. L'inégalité de Heisenberg est dite saturée pour un paquet d'onde gaussien (cas d'égalité).
  
  % \begin{align*}
  %   \psi(x,t) & = \int_{-\infty}^{+\infty} A_0 \exp\left[\, -\frac{(k-k_0)^2}{4\Delta k^2}\right] \e^{ikx- i\omega(k)t}\d k \\ 
  %   & = A_0 \int_{-\infty}^{+\infty} \exp\left[\, -\frac{k^2}{4\Delta k^2} - \frac{kk_0}{2\Delta k^2} - \frac{k_0^2}{4 \Delta k^2} + ikx- i\frac{\hbar k^2}{2m}t\right]\d k. \\ 
  %   & = A_0 e^{ -\frac{k_0^2}{4 \Delta k^2}} \int_{-\infty}^{+\infty} \e^{-ak^2+b k}\d k,
  % \end{align*}
  % avec $a = \frac{1}{4\Delta k^2} + \frac{i\hbar t}{2m}$ et $b = \frac{k_0}{2\Delta k^2} + ix$.
  
  }

  % Déterminer la position moyenne $\moy{x}$ et la largeur spatiale $\Delta x$ du paquet d’onde et vérifier que l’inégalité de Heisenberg $\Delta x \Delta p \geq \frac{\hbar}{2}$ est satisfaite. Quelle est la loi horaire $\moy{x}\!(t)$ de la  trajectoire classique ?

\item Déterminer la constante $A_0$ à partir de la condition de normalisation. 
\sol{
  On a
  $$
  \int_{-\infty}^{+\infty} \rho(x,t=0) \d x = |A_0|^2 4\pi \Delta k^2  \int_{-\infty}^{+\infty} \e^{-2\Delta k^2 x^2} \d x = |A_0|^2 4\pi \Delta k^2 \sqrt{\frac{\pi}{2\Delta k^2}}  = |A_0|^2 \Delta k (2\pi)^{3/2}  = 1,
  $$
  soit 
  \eqbox{
    A_0 = \frac{1}{(2\pi)^{3/4}} \frac{1}{\sqrt{\Delta k}}.
  }
  à une phase globale près. On peut choisir $A_0$ réel et positif sans perte de généralité.
}
% Déterminer la vitesse de phase~$v_\vphi$ et la vitesse de groupe $v_g$. Comment retrouver la vitesse \guill{classique} de la particule ? 

\item Rappeler la relation de dispersion $\omega(k)$ pour une particule libre. Pour une densité spectrale gaussienne $A(k)$ centrée en un vecteur d'onde moyen $k_0$, la position moyenne vérifie $\moy{x} = \frac{\hbar k_0}{m}t$. Commenter.
\sol{
  La relation de dispersion pour une particule libre est $\omega(k) = \frac{\hbar k^2}{2m}$. On a donc
  $$
  \moy{x} = \int_{-\infty}^{+\infty} x |\psi(x,t)|^2 \d x = \frac{\hbar k_0}{m}t.
  $$
  La vitesse de phase est donnée par $v_\vphi = \frac{\omega(k_0)}{k_0} = \frac{\hbar k_0}{2m}$ et la vitesse de groupe par $v_g = \dd{\omega(k)}{k}\big|_{k=k_0} = \frac{\hbar k_0}{m}$, soit 
  \eqbox{
  \moy{x} = v_g t.
  }
  La position moyenne du paquet d'onde suit donc la trajectoire classique de la particule libre.
}

\item On représente ci-dessous la partie réelle de $\psi(x,t)$ à différents instants avec des unités de temps et de distance arbitraires. Commenter ces graphes à la lumière des résultats précédents. Comment varie l'indétermination sur la position $\Delta x(t)$ au cours du temps ?
\sol{
On voit que l'indétermination $\Delta x$ augmente avec le temps, ce qui est dû à la dispersion du paquet d'onde. autrement dit, le paquet d'onde s'étale. Cette propriété est caractéristique des équations de diffusion. Physiquement, cela signifie que la particule se délocalise de plus en plus au cours de sa propagation, ce qui est logique car sa vitesse initiale est aussi indéterminée.
}
\end{quest}
\textsc{Données} : $\int_{-\infty}^{+\infty} \e^{-au^2+bu}\d u = \sqrt{\frac{\pi}{a}}\,\e^{\frac{b^2}{4a}}.$

%\begin{figure}[h!]
%\centering
\begin{center}
% \includegraphics[width=0.48\textwidth]{bouffees_onde.png}
\includegraphics[width=0.8\columnwidth]{paquet.png}
\end{center}
%\end{figure}
\end{exo}
%%%%%%%%%%%%%%%%%%%%%%%%%%%%%%%%%%%%%%%%%%%%%%%%%%%%%%%%%%%%%%%%%%%%%%%%%%%%%%%%%%%

%%%%%%%%%%%%%%%%%%%%%%%%%%%%%%%%%%%%%%%%%%%%%%%%%%%%%%%%%%%%%%%%%%%%%%%%%%%%%%%%%%%
\begin{exo}[PC]{Modèle de l'atome d'hydrogène}
On considère un atome d'hydrogène sphérique de rayon $a$. L'énergie de l'électron dans l'atome se met sous la forme
\begin{equation*}
E= \frac{p^2}{2m}-\frac{e^2}{4\pi\varepsilon_0 a}.
\end{equation*}
\begin{quest}
\item Interpréter physiquement chaque terme de l'énergie.
\item Estimer l’indétermination sur la position $\Delta x$, puis celle sur l’impulsion $\Delta p$. En déduire l’expression de $E$ en fonction du rayon $a$ uniquement.

\item Pour quelle valeur de $a$, notée $a_{0}$, l’énergie $E$ est-elle minimale ? Que vaut  l’énergie minimale associée $E_0$ ? Faire l'application numérique.
% Ce calcul donne l'ordre de grandeur de la taille de l'atome d'hydrogène.
% \item Déterminer la valeur minimale de l'expression approchée de $E$ et la calculer numériquement.
\item
En mécanique classique, un électron en orbite autour du noyau perdrait de l’énergie sous forme de rayonnement.
%électromagnétique et
%, on trouve une énergie mécanique $E=-\frac{e^2}{8\pi\varepsilon_0 a}$. De plus l'électron en mouvement perd de l'énergie par rayonnement électromagnétique.
 Expliquer pourquoi l'inégalité de Heisenberg est responsable de la stabilité des atomes.
\end{quest}
\end{exo}
%%%%%%%%%%%%%%%%%%%%%%%%%%%%%%%%%%%%%%%%%%%%%%%%%%%%%%%%%%%%%%%%%%%%%%%%%%%%%%%%%%%

%%%%%%%%%%%%%%%%%%%%%%%%%%%%%%%%%%%%%%%%%%%%%%%%%%%%%%%%%%%%%%%%%%%%%%%%%%%%%
\begin{exo}[PC]{Interférences d'atomes de néon}
En 1992, Shimizu \emph{et al.} observèrent l'interférence entre atomes de néon passant dans une double-fente d'Young. Les atomes sont d'abord piégés par un système magnéto-optique et refroidis à la température $T=2,5~\mu$K. Le piège est ensuite désactivé, de sorte que les atomes tombent en chute libre à travers une double-fente située à une distance $d=76$ mm sous le piège, puis frappent l'écran placé à une distance $D=113$ mm en dessous des fentes. La largeur des fentes est $b=2~\mu$m et leur espacement $a=6~\mu$m. Un atome de néon a une masse $m=3,35 \cdot 10^{-26}$ kg. La constante de Planck vaut $h=6,63 \cdot 10^{-34}$ 	{J}{}{s}{}
\begin{quest}
	\item On néglige la vitesse d'agitation thermique des atomes. Après avoir estimé le temps de chute $t$, déterminer la longueur de De Broglie $\lambda$ au niveau de l'écran puis calculer l'interfrange $i$ de la figure d'interférence.
	\item Comparer l'interfrange calculée à celle déterminée à partir des courbes expérimentales de la Fig.~\ref{fig:interf}. Discuter de l'écart entre ces deux valeurs.
\end{quest}
%\begin{figure}[h!]
\begin{center}
\includegraphics[width=0.3\columnwidth]{cropped_p03.pdf}
%\caption{Figures d'interférence obtenues par Shimizu \emph{et al.} pour différents temps de chute $t$. La dernière courbe est la prédiction théorique.}
%\label{fig:interf}
\end{center}
\end{exo}
%%%%%%%%%%%%%%%%%%%%%%%%%%%%%%%%%%%%%%%%%%%%%%%%%%%%%%%%%%%%%%%%%%%%%%%%%%%%%

